% Standalone abstract.  Compile: pdflatex abstract.tex
\documentclass[11pt,a4paper]{article}
\usepackage[margin=25mm]{geometry}
\usepackage[T1]{fontenc}\usepackage{lmodern}\usepackage{amsmath,amssymb,mathtools}
\usepackage{siunitx}
\sisetup{detect-all=true,per-mode=symbol,exponent-product=\times,group-digits=integer,group-minimum-digits=5,output-decimal-marker={.},group-separator={\,}}
\newcommand{\vect}[1]{\boldsymbol{#1}}\newcommand{\dd}{\mathrm{d}}\newcommand{\ii}{\mathrm{i}}\newcommand{\ee}{\mathrm{e}}
\DeclareMathOperator{\Real}{Re}\DeclareMathOperator{\Imag}{Im}
\pagestyle{empty}
\begin{document}
\begin{center}\bfseries \bfseries Backward-branch negative reflection of guided waves from the
free edge of a flexoelectric piezoelectric semiconductor plate\end{center}
\medskip
\noindent\textbf{Abstract}\par\smallskip
\noindent
This paper presents a model of the scattering of a guided wave at the free edge of an n-type ZnO plate ($6mm$ symmetry), in which strain-gradient elasticity, piezoelectricity, flexoelectricity and linearized electron drift--diffusion act together, with a carrier--strain-gradient coupling added to the electric enthalpy.  The governing equations follow from a variational principle, which also yields the edge conditions, including the double traction and the corner line force.  For a wave travelling along the plate they reduce to a quartic eigenvalue problem, which is solved at fixed frequency by a spectral Galerkin expansion through the thickness; the incident mode is then matched to the outgoing modes at the traction-free edge.  The partition of the incident power shows that a backward mode, whose energy flux is opposite to its phase velocity, carries \SI{1.93}{\percent} of the power at the chosen operating point and up to \SI{94.29}{\percent} near a zero-group-velocity point of the same branch.  These are modal shares and not reflectances, because the incident flux itself becomes small there.  The conversion belongs to the backward branch of the plate and not to the added coupling, which acts mainly through conduction loss and lowers the share, and near the zero-group-velocity point the share is sensitive to carrier density and to elastic loss.

\medskip
\noindent\textbf{Keywords:} flexoelectric semiconductor; strain-gradient
elasticity; guided plate waves; backward branch; negative reflection; spectral
Galerkin method.
\end{document}
